%\usepackage [ISO-8859-1]{inputenc}
\input{Olympiad}

\thispagestyle{fancy}
\lstset{basicstyle=\ttfamily,
  breaklines=true}
% ----------------------------------------------------

% Some ugly hacking to ensure the intro stays 1 page long
\newlength{\oldparskip}
\setlength{\oldparskip}{\parskip}

\setlength{\parskip}{0.5pt}

\begin{center}
\Huge{Programmeringsolympiaden 2026}
\end{center}
\vspace{-1.2cm} 
\specialsection*{Tävlingsregler för skolkvalet}
{
\setlength{\itemsep}{0.5pt}
\begin{lista}
\item Tävlingen äger rum på av skolan bestämt datum under {\bf fyra
    timmar. Ingen förlängning ges för lunch eller raster.} Eleven ska i förväg komma överens med läraren om att använda egen dator eller en som skolan tillhandahåller. I vilket fall som helst måste eleven befinna sig i avtalad lokal på skolan. 

\item Tävlingen består av sex uppgifter som vardera ska lösas genom ett
  datorprogram i valfritt programmeringsspråk.
\item {\bf Indata kan läsas in i programmet på valfritt sätt}, t.ex. genom att programmet för en dialog med användaren (som i körningsexemplena i uppgifterna), att de skrivs in i ett grafiskt gränssnitt eller att datafiler skickas till {\em standard input}. Kom bara överens med din lärare om hur programmet ska testas.
\item Dina lösningar kommer att testköras med förpreparerade
  indata. Varje uppgift testas normalt med 5 testfall, som vardera ger
  1 poäng om ditt program skriver ut korrekt svar inom en exekveringstidsgräns av
  {\bf 3 sekunder}. Ingen test av indata behöver göras, den följer specifikationerna 
i uppgiften.
\item Det är ofta olika begränsningar på de olika testfallen, t.ex.\ storleken på
indata eller andra inskränkningar. Detta anges i uppgiften. {\bf Observera att det kan vara
helt olika svårighetsgrad på en uppgift beroende på dessa skillnader. Det kan därför vara
lättare att få delpoäng på en uppgift som verkar svår än att få full
poäng på en uppgift som verkar lättare.} Informationen om delpoäng är
därför extremt viktig för att planera sin
tävling.
\item Rättningen utförs på samma eller likvärdig dator. Ändringar
  i källkoden tillåts ej efter tävlingen. Om programmet
  inte kan kompileras ges 0 p.\ på uppgiften. 
\item Om något av följande inträffar ger
  det {\em testfallet} 0 poäng, men programmet fortsätter testas med övriga testfall.
\begin{lista}
\item Exekveringstiden överstiger 3 sekunder \vspace{-0.2cm}
\item Exekveringsfel (run time error) \vspace{-0.2cm}
\item Fel svar \vspace{-0.2cm}
\end{lista}

\item Deltagandet är individuellt vilket bland annat innebär att inget utbyte av idéer eller 
filer får ske under tävlingen.  
\item Hjälpmedel: Valfritt skriftligt material, material som finns installerat på datorn samt material som finns tillgängligt på internet.
Det är \ku{inte} tillåtet att aktivt kommunicera på internet (t.ex. chatta eller ställa frågor till ett forum) utan endast att söka efter information.
{\bf Det är heller inte tillåtet att använda generativ AI såsom ChatGPT eller Github Copilot.}
Miniräknare är tillåten.
\item Tävlingsbidraget ska lämnas in i form av källkodsfiler (uppg1...uppg6 med passande filtillägg) som läggs på överenskommen plats. Ingen hänsyn tas till andra filer. Var noga 
med att lämna in den korrekta versionen av ditt program.  
\end{lista}
}
De högst placerade i kvalet går vidare till finalen i Göteborg där landslagsplatser till BOI i Finland och IOI i Uzbekistan står på spel. Lycka till!
\begin{center}
% \Large Lycka till!
\end{center}
%\end{document}

\setlength{\parskip}{\oldparskip}

\newpage
\specialsection*{Uppgift 1 -- Triangelfabriken}

Alice och Bertil kastar boll i en park. Exakt $21$ meter framför Alice och Bertil finns en stolpe, och Jack utmanar Alice och Bertil till en liten utmaning.

Utmaningen är att Alice och Bertil ska kasta varsin boll så nära stolpen som möjligt, men bollen får inte kastas förbi stolpen.
Den som kastar närmast $21$ meter, utan att kasta förbi $21$ meter, vinner. Om både Alice och Bertil kastar över $21$ meter, eller om de kastar lika långt så vinner Jack. De har nu kastat och hjälp att avgöra vem som vann.

Skriv ett program som läser in talen $A$ och $B$, antalet meter som Alice och
bertil har kastat vardera, och sedan skriver ut vem som vunnit.

Det är garanterat att $(0 \leq A,B \leq 30)$ håller.

Sedan ska programmet skriva vem som vann: antingen ''Alice'', ''Bertil'' eller ''Jack''. Notera att du inte ska skriva ut citationstecken.

\fe{Poängsättning:}\\
De fem testfallen som vardera ger $1$ poäng har följande begränsningar:
\begin{itemize}
  \item[(1)] $\;\;\;\;$ $A \neq B$ och $A,B \leq 21$
  \item[(2)] $\;\;\;\;$ $A, B \leq 21$
  \item[(3)] $\;\;\;\;$ Jack vann inte i det här testfallet.
  \item[(4-5)] $\;\;\;\; $ Inga ytterligare begränsningar.
\end{itemize}


\vspace{1cm}
\fe{Körningsexempel 1}
\begin{verbatim}
Alice kast ? 7
Bertils kast ? 5

Svar: Alice
\end{verbatim}
\fe{Förklaring:} Alice är närmare $21$, därmed vann hon.

\vspace{1cm}

\fe{Körningsexempel 2}
\begin{verbatim}
Alice kast ? 26
Bertils kast ? 22

Svar: Jack
\end{verbatim}
\fe{Förklaring:} Eftersom båda kastade längre än $21$ meter vann Jack.

\newpage
\specialsection*{Uppgift 2 -- Snurra kompass}


%\begin{center}
%  \includegraphics[width=12cm]{sifferkrypto.pdf} \\ ~\\
%  \ku{Exempel på missuppfattning som kan ske med Alice nya chiffer.}
%\end{center}

Du har en kompass som kan peka i $360$ olika riktningar. Just nu pekar den rakt uppåt, på talet $360$.

Tyvärr så funkar inte kompassen, och det enda sättet du kan få den att peka i andra riktningar
är genom lyfta på hela kompassen och rotera den ett antal grader motsols. Till din hjälp har du
två robotarmar. Dessa är tyvärr inte perfekta: den första robotarmen kan
snurra kompassen exakt $a$ grader motsols och den andra robotarmen kan snurra $b$ grader motsols.
Robotarmarna kan endast rotera kompassen $a$ respektive $b$ grader varje gång
de används, varken mer eller mindre.
På grund av detta går det kanske inte att få kompassen att peka
i alla $360$ riktningar.

Detta får dig att undra hur många unika riktningar kompassen kan peka åt?

Skriv ett program som läser in talen $a$ och $b$, antalet grader som den första respektive andra robotarmen snurrar kompassen.

Det är garanterat att $(1 \leq a,b \leq 360)$ håller.

Sedan ska programmet skriva ut antal unika riktningar som kompassen kan peka åt om man snurrar den ett flertal gånger med robotarmarna.

\fe{Poängsättning:}\\
De fem testfallen som vardera ger $1$ poäng har följande begränsningar:
\begin{itemize}
  \item[(1)] $\;\;\;\;$ $a = 13$
  \item[(1-2)] $\;\;\;\;$ $a = b$
  \item[(3-5)] $\;\;\;\; $ Inga ytterligare begränsningar.
\end{itemize}

\vspace{0.3cm}

\fe{Körningsexempel 1}
\begin{verbatim}
Gradtal a ? 50
Gradtal b ? 180

Svar: 36
\end{verbatim}
\fe{Förklaring:} Med denna uppsättningen robotarmar går det att få kompassen
att peka på varje heltalsmultipel av $10$: $0, 10, 20, \dots, 350$.
Notera att $0$ och $360$ är samma gradtal. Om hon använder robotatarm $a$
$8$ gånger kan hon exempelvis på kompassen att peka mot $40$ grader.

\vspace{0.7cm}

\fe{Körningsexempel 2}
\begin{verbatim}
Gradtal a ? 360
Gradtal b ? 360

Svar: 1
\end{verbatim}
\fe{Förklaring:} Hur många gånger du än roterar kompassen kommer den fortsätta vara riktad åt $0$.


\newpage
\specialsection*{Uppgift 3 -- Arborist}

Teo har äntligen landat sitt drömjobb; Han har börjat jobba som arborist!
En viktig del av en arborists arbete är att plantera träd.

Teo behöver nu hjälp med att plantera $N$ träd längs med en tallinje.
För alla $N$ träd, har Teo redan bestämt det koordinat längs tallinjen som trädet ska planteras vid. 
De bestämda kooridnaterna är garanterat att vara positiva heltal.
Teo befinner sig just nu vid koordinat $0$, tillsammans med $N$ stora trädplantor.
För att plantera träden kan han göra följande:
\begin{itemize}
  \item Röra sig ett steg åt höger. Detta tar en sekund.
  \item Röra sig ett steg åt vänster. Detta tar en sekund.
  \item Om han befinner sig vid koordinat $0$ kan han plocka upp ett av de $N$ trädplantorna. Detta tar en sekund.
  \item Ta ett av trädplantorna han bär på, och plantera det där han befinner sig. Detta tar en sekund.
\end{itemize}
Tyvärr orkar Teo bära som mest $K$ kilogram träd åt gången. Varje träd väger antingen $20$ eller $40$ kilogram.

När han väl planterat alla träd måste han återvända till koordinat $0$. Han behöver nu din hjälp att lista ut hur lång tid detta kommer ta om han gör detta på snabbaste möjligaste vis.

Skriv ett program som först läser in $N$ och $K$ ($1 \leq N \leq 8$, $20 \leq K \leq 50$), antalet träd, och antalet kilogram som Teo orkar bära.
På nästa rad ska  programmet därefter läsa in de $N$ koordinaterna $x_1, \dots, x_N$ ($1 \leq x_i \leq 30$) där varje trädplanta ska planteras. På sista raden ska det läsa in de $N$ vikterna $k_1, \dots, k_N$ ($k_i \leq K$),  hur mycket varje trädplanta väger. Det är även garanterat att varje träd antingen väger $20$ eller $40$ kilogram.

Programmet ska sedan skriva ut antalet sekunder det tar för Teo att plantera alla träd och återvända till $0$ om han gör detta på det snabbaste möjliga sättet.

\vspace{0.5cm}
\fe{Poängsättning:}\\
De fem testfallen som vardera ger $1$ poäng har följande begränsningar:
\begin{itemize}
    \item[(1)] $\;\;\;\;$ $K=25$
    \item[(2)] $\;\;\;\;$ $k_i = 20$
    \item[(3)] $\;\;\;\;$ Alla träd ska planteras på samma $x$-koordinat.
    \item[(4-5)] $\;\;\;\; $ Inga ytterligare begränsningar.
\end{itemize}

\vspace{0.5cm}
\fe{Körningsexempel 1}
\begin{verbatim}
N ? 2
K ? 45
Trädens koordinater ? 1 2
Trädens vikt ? 20 20

Svar: 8
\end{verbatim}
\fe{Förklaring:} Denna indatan betyder att det finns två träd. Ett väger $20$kg och ska planteras vid koordinat $1$, och det andra väger också $20$kg och ska planteras vid koordinat $2$.
I exempelfallet är det optimalt att för Teo att göra följande:
\begin{itemize}
  \item Plocka upp trädet med koordinat $1$.
  \item Plocka upp trädet med koordinat $2$.
  \item Gå höger.
  \item Plantera trädet med koordinat $1$.
  \item Gå höger.
  \item Plantera trädet med koordinat $2$.
  \item Gå vänster.
  \item Gå vänster. Han har nu återvänt till koordinat $0$ och är därmed färdig.
\end{itemize}
Det finns andra sätt att plantera träden lika snabbt, men det går inte på
mindre än $8$ sekunder.




\newpage
\specialsection*{Uppgift 4 -- Tingsplatsen}

Arkeologen Arvid är på jakt efter en okänd tingsplats från järnåldern någonstans i Mellansverige. Han har kartlagt
ett antal gårdar som fanns i området, och har kommit fram till att tingsplatsen måste ligga på samma avstånd
till de olika gårdarna.

Området som Arvid undersökt kan representeras av ett rutnät med $N$ rader och $M$ kolumner. Det tar en minut för en
järnåldersmänniska att ta sig mellan två närliggande rutor. Två rutor är närliggande om de delar en av sina fyra sidor.

Du får givet en karta av rutnätet som markerar alla positioner av gårdar. 
Det är känt att det tar lika lång tid att ta sig från var och en av gårdarna till den okända tingsplatsen. 

Din uppgift är att skriva ett program som hittar vilka rutor där tingsplatsen skulle kunna funnits.

Programmet ska först läsa in heltalen $N$ och $M$ ($1 \leq N,M \leq 10$), antalet rader och kolumner i rutnätet. Läs sedan $N$ rader, som innehåller
rutnätets rader.
Tecknet ''\verb|*|'' betyder att det finns en gård på motsvarande plats i rutnätet, och
''\verb|.|'' betyder att det inte finns en gård.

Programmet ska sedan skriva ut en kopia av rutnätet i indatan, men där det markerat med ''\verb|X|'' de rutor som tingsplatsen kan vara belägen på.

\vspace{0.5cm}
\fe{Poängsättning:}\\
De fem testfallen som vardera ger $1$ poäng har följande begränsningar utöver de som redan nämnts:
\begin{itemize}
    \item[(1)] $\;\;\;\;$ $N = 1$
    \item[(2)] $\;\;\;\;$ Det finns exakt $4$ gårdar som alla ligger på hörnet av en rektangel.
    \item[(3-5)] $\;\;\;\;$ Inga ytterligare begränsningar.
\end{itemize}

\vspace{0.5cm}
\fe{Körningsexempel 1}\\
\vspace{-0.3cm}
\begin{verbatim}
N ? 4
M ? 5
Rad 1 ? *...*
Rad 2 ? .....
Rad 3 ? .....
Rad 4 ? .....

Svar:
*.X.*
..X..
..X..
..X..

\end{verbatim}
\fe{Förklaring:} TODO


\specialsection*{Uppgift 5 -- Miniräknaren}


Det sägs att uppe i Alperna bor det $C$ stycken systrar som vardera har ett eget favorittal $f$.
Varje syster har därför byggt ett hus $f$ höjdenheter över vattenytan.
Dessa husen är byggda på ett koordinatsystem.
Den första systern har byggt sitt hus vid $(0,0)$, den andra vid $(0,1)$ och så vidare.

När husen befinner sig så högt uppe blir det svårt att ta sig upp och ner.
Därför har varje syster byggt en trappa genom att bygga uppåt, där hon
sänkar höjden med $1$ för varje koordinat hon går uppåt.
Exempelvis, om systerns hus befinner sig på $(0,2)$ och har höjd $10$, så bygger hon ett trappsteg vid $(1,2)$ med höjd $9$, ett trappsteg vid $(2,2)$ med höjd $8$, och så vidare.

Du befinner dig just nu vid dessa trappor.
Mer exakt befinner du dig på koordinaterna $(s_x, s_y)$ och vill ta dig till $(g_x, g_y)$.
Du kan röra dig till rutor som intilliggande antingen upp, vänster, ner eller höger.
Det är aldrig problem att gå höger eller vänster, även om det finns stora höjdskillnader.
Självfallet kan du inte röra dig till en koordinat som varken har ett hus eller ett trappsteg.
Du vägrar även att röra dig till en koordinat vars $y$-värde är större än eller lika med talet $R$.

De höga höjderna är ganska jobbiga för dig, så ditt mål är att minimera tiden du spenderar på hög höjd.
Mer exakt, om du under din vandring från $(s_x, s_y)$ till $(g_x, g_y)$ besöker en viss mängd rutor, så definierar vi jobbigheten av vandringen som summan av höjden av alla rutor du besöker.
Om du vandrar optimalt, vad är minsta möjliga jobbigheten?

\begin{figure}[h]
  \centering
  \includegraphics[width=0.3\textwidth]{sample1.pdf}
    \\Bilden visar hur höjderna ser ut i exempelfall 1. Raden längst ner innehåller systrarnas hus.
    Start- och mål är även markerade, samt den mest minst jobbiga vandringar mellan dem.
\end{figure}

Indataformat bla bla jag orkar inte
Skriv ett program som först läser in heltalet $N$ ($1 \le N \le 10^9$), som utgör kålnami-koden.
Det ska sedan läsa in siffran $M$ ($2 \leq M \leq 9$), som beskrivits ovan.

Programmet ska sedan skriva ut jobbigheten av den minst jobbiga vandring från start till koordinaterna du vill ta dig till.
\textbf{Notera: svaret ryms inte nödvändigtvis i ett 32-bitars heltal.}

\fe{Poängsättning:}\\
De fem testfallen som vardera ger $1$ poäng har följande begränsningar utöver de som redan nämnts:
\begin{itemize}
    \item[(1)] $\;\;\;\;$ $C=2$ 
    \item[(2-3)] $\;\;\;\;$ $R \leq 10^4$ 
    \item[(4-5)] $\;\;\;\;$ Inga ytterligare begränsningar.
\end{itemize}


\fe{Körningsexempel 1}
\begin{verbatim}
R ? 4
C ? 2
Bottenraden ? 15 12 5
Start-x ? 0
Start-y ? 0
Mål-x ? 0
Mål-y ? 4

Svar: 61
\end{verbatim}

\fe{Förklaring:}
TODO

Se nästa sida för ett till körningsexempel.

\vspace{1cm}

\fe{Körningsexempel 2}
\begin{verbatim}
R ? 5
C ? 3
Bottenraden ? 5 20 25
Start-x ? 1
Start-y ? 4
Mål-x ? 2
Mål-y ? 0

Svar: 76
\end{verbatim}

\vspace{1cm}

\fe{Körningsexempel 3}
\begin{verbatim}
R ? 100000
C ? 6
Bottenraden ? 218849 99999 100019 100003 100000 264319
Start-x ? 1
Start-y ? 8493
Mål-x ? 5
Mål-y ? 99998

Svar: 4186884119
\end{verbatim}


\newpage
\specialsection*{Uppgift 6 -- Binära palindrom}


Mikael studerar intressanta positiva heltal, och vill nu undersöka tal som är ett \emph{binärt palindrom}. 
För att formellt beskriva detta definierar vi följande funktioner:

\[
\text{bin}(n) = \text{den binära representationen av heltalet } n \text{ utan inledande nollor},
\]

\[
\text{reverse}(s) = \text{strängen } s \text{ skriven baklänges}.
\]

Ett heltal $n$ är alltså ett \emph{binärt palindrom} om
\[
\text{bin}(n) = \text{reverse}(\text{bin}(n)).
\]

Notera att ett palindrom är en sträng som läses likadant framifrån som bakifrån.


Mikael gillar de binära palindromer som innehåller exakt $k$ ettor i dess binära representation, 
och vill nu veta hur många sådana tal som finns inom intervallet $[l,r]$.

Givet tre heltal $l$, $r$ och $k$, bestäm hur många heltal $x$ som uppfyller:
\begin{enumerate}
    \item Talet $x$ är mellan $l$ och $r$, alltså att $l \leq x \leq r$,
    \item Talet $x$ är ett tal som Mikael gillar (dvs.\ ett binärt palindrom),
    \item Binärrepresentationen av $x$ innehåller exakt $k$ ettor.
\end{enumerate}

Skriv ett program som läser in heltalen $l, r, k$ $l, r, k$ ($1 \leq l \leq r \leq 10^{18}, 1 \leq k < 60$).

Programmet ska skriva ut antalet tal i intervallet $[l, r]$ som är binära palindromer och har exakt $k$ ettor i sin binära representation.

\fe{Poängsättning:}\\
De fem testfallen som vardera ger $1$ poäng har följande begränsningar utöver de som redan nämnts:
\begin{itemize}
    \item[(1)] $\;\;\;\;$ $l,r \leq 10^9$
    \item[(2)] $\;\;\;\;$ $k \leq 4$
    \item[(3-5)] $\;\;\;\; $ Inga ytterligare begränsningar.
\end{itemize}


\fe{Körningsexempel 1}
\begin{verbatim}
l, r, k ? 7 10 2

Svar: 1
\end{verbatim}
\fe{Förklaring:} 
Vi söker tal i intervallet $[7, 10]$ som både är binära palindromer och har exakt två ettor i sin binära representation.

Det finns endast 4 tal i intervallet:

\begin{align*}
7 &= 111_2 \quad (\text{palindrom, tre ettor}) \\
8 &= 1000_2 \quad (\text{ej palindrom}) \\
9 &= 1001_2 \quad (\text{palindrom, två ettor}) \\
10 &= 1010_2 \quad (\text{ej palindrom})
\end{align*}

Det enda talet som uppfyller båda villkoren är $9$, så svaret är $1$.

\vspace{1cm}

\fe{Körningsexempel 2}
\begin{verbatim}
l, r, k ? 1 6 2

Svar: 2
\end{verbatim}
\fe{Förklaring:} 
Vi söker tal i intervallet $[1, 6]$ som både är binära palindromer och har exakt två ettor i sin binära representation.

Det finns endast 6 tal i intervallet:

\begin{align*}
1 &= 1_2 \quad (\text{palindrom, en etta}) \\
2 &= 10_2 \quad (\text{ej palindrom}) \\
3 &= 11_2 \quad (\text{palindrom, två ettor}) \\
4 &= 100_2 \quad (\text{ej palindrom}) \\
5 &= 101_2 \quad (\text{palindrom, två ettor}) \\
6 &= 110_2 \quad (\text{ej palindrom})
\end{align*}

Talen som uppfyller båda villkoren är $3$ och $5$, så svaret är $2$.

\vspace{1cm}

\fe{Körningsexempel 3}
\begin{verbatim}
l, r, k ? 999 99999 6

Svar: 94
\end{verbatim}

\end{document}
